% !TeX program = xelatex
\documentclass[a4paper,14pt]{extarticle}

\usepackage{pdfpages}

% Язык и основной шрифт.
\usepackage{fontspec}
\setmainfont{Times New Roman}

\usepackage[main=russian,english]{babel}
\emergencystretch=3em % растягивать пробелы при overfull
% Поля страницы по требованиям ИнЭУ.
\usepackage[
  left=30mm,
  right=10mm,
  top=20mm,
  bottom=20mm,
  bindingoffset=0cm
]{geometry}


% Абзац, интервал и выравнивание текста.
\usepackage{setspace}
\onehalfspacing
\setlength{\parindent}{1.25cm}
\setlength{\parskip}{0pt}
\usepackage{indentfirst}
\usepackage{ragged2e}
\usepackage{microtype}
\justifying

% Математика и формулы.
\usepackage{amsmath}
\usepackage[
  math-style=TeX,
  bold-style=TeX
]{unicode-math}

\setmathfont{TeX Gyre Termes Math}

% Рисунки и таблицы.
\usepackage{graphicx}
\graphicspath{{fig/}}
\usepackage{array,tabularx,multirow}
\usepackage{float}
\usepackage[table]{xcolor}


% Единая нумерация рисунков, таблиц и формул внутри раздела.
\numberwithin{figure}{section}
\numberwithin{table}{section}
\numberwithin{equation}{section}
% Для сквозной нумерации закомментировать три строки выше.

% Подписи рисунков и таблиц.
\usepackage{caption}
\DeclareCaptionLabelSeparator{vkrdash}{\space--\space}
\captionsetup[figure]{
  name=Рисунок,
  labelsep=vkrdash,
  justification=centering,
  singlelinecheck=false,
  position=bottom,
  skip=6pt
}
\captionsetup[table]{
  name=Таблица,
  labelsep=vkrdash,
  justification=raggedright,
  singlelinecheck=false,
  position=top,
  skip=6pt
}

% Заголовки разделов и подразделов.
\usepackage{titlesec}
\titleformat{\section}[block]
  {\normalfont\normalsize\bfseries\raggedright}
  {\thesection}{1em}{}
\titleformat{name=\section,numberless}[block]
  {\normalfont\normalsize\bfseries\centering}
  {}{0pt}{\MakeUppercase}
\titleformat{\subsection}[block]
  {\normalfont\normalsize\bfseries\raggedright}
  {\thesubsection}{1em}{}
\titleformat{\subsubsection}[block]
  {\normalfont\normalsize\bfseries\raggedright}
  {\thesubsubsection}{1em}{}


% Интервалы до и после заголовков.
\titlespacing*{\section}{\parindent}{\baselineskip}{0pt}
\titlespacing*{name=\section,numberless}{0pt}{\baselineskip}{0pt}
\titlespacing*{\subsection}{\parindent}{\baselineskip}{0pt}
\titlespacing*{\subsubsection}{\parindent}{\baselineskip}{0pt}


% Содержание.
\usepackage{tocloft}
\renewcommand{\contentsname}{СОДЕРЖАНИЕ}
\addto\captionsrussian{\renewcommand{\contentsname}{СОДЕРЖАНИЕ}}
\renewcommand{\cfttoctitlefont}{\hfill\normalfont\normalsize\bfseries}
\renewcommand{\cftaftertoctitle}{\hfill}
\setlength{\cftbeforetoctitleskip}{0pt}
\setlength{\cftaftertoctitleskip}{\baselineskip}

\renewcommand{\cftsecfont}{\normalfont}
\renewcommand{\cftsecpagefont}{\normalfont}
\renewcommand{\cftsubsecfont}{\normalfont}
\renewcommand{\cftsubsecpagefont}{\normalfont}

\renewcommand{\cftsecleader}{\cftdotfill{\cftdotsep}}
\renewcommand{\cftsubsecleader}{\cftdotfill{\cftdotsep}}
\renewcommand{\cftsubsubsecleader}{\cftdotfill{\cftdotsep}}

% В содержание выводятся только section и subsection; subsubsection скрыты.
\setcounter{tocdepth}{2}

% Отступы в содержании: разделы без отступа, подразделы с отступом.
\setlength{\cftbeforesecskip}{0pt}
\setlength{\cftbeforesubsecskip}{0pt}
\setlength{\cftsecindent}{0pt}
\setlength{\cftsubsecindent}{2em}
\setlength{\cftsubsubsecindent}{4em}


% Списки.
\usepackage{enumitem}
\makeatletter
\AddEnumerateCounter*{\asbuk}{\c@asbuk}{7}
\makeatother
\setlist[itemize]{
  label=--,
  leftmargin=1.25cm,
  itemsep=0pt,
  topsep=0pt
}
\setlist[enumerate,1]{
  label=\arabic*),
  leftmargin=1.25cm,
  itemsep=0pt,
  topsep=0pt
}
\setlist[enumerate,2]{
  label=\asbuk*),
  leftmargin=1.25cm,
  itemsep=0pt,
  topsep=0pt
}

% Библиография по ГОСТ-стилю biblatex.
\usepackage{booktabs}
\usepackage{csquotes}
\usepackage[
  backend=biber,
  style=gost-numeric,
  sorting=none,
  autolang=other,
  language=russian,
  giveninits=true,
  maxnames=3,
  minnames=3,
  doi=false,
  isbn=false,
  url=true,
  eprint=false,
  alldates=short
]{biblatex}
\addbibresource{src/references.bib}

% Чтобы вместо «дата обр.» печаталось «дата обращения: ...»
\DefineBibliographyStrings{russian}{
  urlseen = {дата обращения}
}
\DeclareFieldFormat{urldate}{%
  \mkbibparens{\bibstring{urlseen}\addcolon\space#1}%
}

% Убрать служебные поля из вывода списка.
\AtEveryBibitem{
  \clearfield{note}
  \clearfield{abstract}
  \clearfield{keywords}
  \clearfield{isbn}{}
}

% =========================================================
% Единый стиль библиографии: без курсива, жирного и т. д.
% =========================================================

% ----- 1. Глобальные обёртки biblatex -----
% mkbibemph  — курсив (используется gost-numeric для имён и ряда полей)
% mkbibbold  — полужирный (тома, номера конференций в некоторых стилях)
% mkbibacro  — акронимы (DOI, URL, ISBN — обычно капитель)
\renewcommand*{\mkbibemph}[1]{#1}
\renewcommand*{\mkbibbold}[1]{#1}
\renewcommand*{\mkbibacro}[1]{#1}

% ----- 2. Части имён авторов -----
\renewcommand*{\mkbibnamefamily}[1]{#1}
\renewcommand*{\mkbibnamegiven}[1]{#1}
\renewcommand*{\mkbibnameprefix}[1]{#1}
\renewcommand*{\mkbibnamesuffix}[1]{#1}

% ----- 3. Названия (заглавия) -----
% Глобально — для типов, не перечисленных ниже
\DeclareFieldFormat{title}{#1}
\DeclareFieldFormat{booktitle}{#1}
\DeclareFieldFormat{maintitle}{#1}
\DeclareFieldFormat{journaltitle}{#1}
\DeclareFieldFormat{issuetitle}{#1}
\DeclareFieldFormat{series}{#1}
\DeclareFieldFormat{eventtitle}{#1}
% По типам записи — biblatex проверяет их раньше глобальных
\DeclareFieldFormat[article]{title}{#1}
\DeclareFieldFormat[book]{title}{#1}
\DeclareFieldFormat[inbook]{title}{#1}
\DeclareFieldFormat[incollection]{title}{#1}
\DeclareFieldFormat[inproceedings]{title}{#1}
\DeclareFieldFormat[online]{title}{#1}
\DeclareFieldFormat[report]{title}{#1}
\DeclareFieldFormat[thesis]{title}{#1}
\DeclareFieldFormat[manual]{title}{#1}
\DeclareFieldFormat[misc]{title}{#1}
\DeclareFieldFormat[patent]{title}{#1}

% ----- 4. Числовые и служебные поля -----
% volume/number иногда выделяются жирным в gost-numeric
\DeclareFieldFormat{volume}{#1}
\DeclareFieldFormat{number}{#1}
% edition: «2-е изд.» иногда курсивом
\DeclareFieldFormat{edition}{#1}
% pages: в некоторых версиях стиля — полужирный диапазон страниц
\DeclareFieldFormat{pages}{#1}

% ----- 5. Метка номера источника в списке — [1], [2] … -----
\DeclareFieldFormat{labelnumberwidth}{#1\adddot}

% ----- 6. URL и дата обращения -----
% Без моноширинного шрифта
\urlstyle{same}
\renewcommand*{\UrlFont}{\normalfont}
% «дата обращения: ДД.ММ.ГГГГ» в скобках
\DefineBibliographyStrings{russian}{urlseen = {дата обращения}}
\DeclareFieldFormat{urldate}{(\bibstring{urlseen}\addcolon\space#1)}

% ----- 7. Принудительный сброс шрифта для всего блока -----
\AtBeginBibliography{
  \normalfont
  \renewcommand{\emph}[1]{#1}%   % ← вот это

  }

% ----- 8. Убрать служебные поля из вывода -----
\AtEveryBibitem{
  \clearfield{note}
  \clearfield{abstract}
  \clearfield{keywords}
  \ifentrytype{article}{\clearfield{isbn}}{}
  \ifentrytype{inproceedings}{\clearfield{isbn}}{}
}

% ----- 9. Межстрочные интервалы в списке -----
\setlength{\bibitemsep}{0pt}
\setlength{\bibparsep}{0pt}

% Ссылки и переносимые URL.
\usepackage{xurl}
% \usepackage[unicode]{hyperref}
\usepackage[unicode,hidelinks]{hyperref} % - для ссылок без цветных рамоук

% Нумерация страниц внизу по центру.
\usepackage{fancyhdr}
\pagestyle{fancy}
\fancyhf{}
\fancyfoot[C]{\thepage}
\renewcommand{\headrulewidth}{0pt}
\renewcommand{\footrulewidth}{0pt}
\fancypagestyle{plain}{%
  \fancyhf{}%
  \fancyfoot[C]{\thepage}%
  \renewcommand{\headrulewidth}{0pt}%
}

%-----------------------------------------------------------

% Структурные элементы ВКР без номера и с записью в содержание.
\newcommand{\structsection}[1]{%
  \clearpage
  \phantomsection
  \section*{#1}%
  \addcontentsline{toc}{section}{\MakeUppercase{#1}}%
}
% Использовать так: \structsection{Введение}, \structsection{Заключение}.

%-----------------------------------------------------------

% Приложения: центрированный заголовок и отдельная нумерация объектов.
\newcommand{\startappendices}{%
  \appendix
  \renewcommand{\thesection}{\Asbuk{section}}%
  \renewcommand{\thefigure}{\thesection.\arabic{figure}}%
  \renewcommand{\thetable}{\thesection.\arabic{table}}%
  \renewcommand{\theequation}{\thesection.\arabic{equation}}%
}
\newcommand{\appendixsection}[1]{%
  \clearpage
  \refstepcounter{section}%
  \setcounter{figure}{0}%
  \setcounter{table}{0}%
  \setcounter{equation}{0}%
  \phantomsection
  \addcontentsline{toc}{section}{ПРИЛОЖЕНИЕ \thesection}%
  \begin{center}
    \textbf{ПРИЛОЖЕНИЕ \thesection}\\[2\baselineskip]
    \textbf{#1}
  \end{center}
}
% Использовать так: \startappendices \appendixsection{Название приложения}.


% Заголовок библиографии как структурный элемент.
\defbibheading{vkrbib}{%
  \clearpage
  \phantomsection
  \section*{СПИСОК ИСПОЛЬЗОВАННЫХ ИСТОЧНИКОВ}%
  \addcontentsline{toc}{section}{СПИСОК ИСПОЛЬЗОВАННЫХ ИСТОЧНИКОВ}%
}
% Печатать список так: \printbibliography[heading=vkrbib]



% Таблицы по ГОСТ: перенос, ячейки, примечания, альбомные страницы, числовые столбцы.
\usepackage{longtable}
\usepackage{makecell}
\usepackage{etoolbox}
\usepackage{threeparttable}
\usepackage{pdflscape}
\usepackage{siunitx}

%% Подпись: «Таблица 1 -- Название» слева над таблицей, без точки в конце.
\DeclareCaptionLabelSeparator{gostdash}{ -- }
\captionsetup[table]{
    name=Таблица,
    labelsep=gostdash,
    justification=raggedright,
    singlelinecheck=false,
    position=top,
    skip=6pt,
    font={normalsize,stretch=1}
}
\captionsetup[longtable]{
    name=Таблица,
    labelsep=gostdash,
    justification=raggedright,
    singlelinecheck=false,
    position=top,
    skip=6pt,
    font={normalsize,stretch=1}
}

%% Текст в таблицах: одинарный интервал, кегль меньше основного, но не менее 10 pt.
%% В классе 14 pt команда \small обычно дает около 12 pt.
\AtBeginEnvironment{tabular}{\small\singlespacing\renewcommand{\arraystretch}{1}}
\AtBeginEnvironment{tabularx}{\small\singlespacing\renewcommand{\arraystretch}{1}}
\AtBeginEnvironment{longtable}{\small\singlespacing\renewcommand{\arraystretch}{1}}

%% Отступы и переносимые таблицы.
\setlength{\tabcolsep}{4pt}
\setlength{\LTpre}{0pt}
\setlength{\LTpost}{0pt}

%% Ячейки и заголовки: центр для заголовков граф, обычное начертание, без точек в конце.
\renewcommand\theadfont{\normalfont}
\renewcommand\theadalign{cc}
\renewcommand\cellalign{tl}
\setcellgapes{2pt}
\makegapedcells

%% Типы столбцов: L — влево, C — по центру, R — вправо.
\newcolumntype{L}[1]{>{\raggedright\arraybackslash}p{#1}}
\newcolumntype{C}[1]{>{\centering\arraybackslash}p{#1}}
\newcolumntype{R}[1]{>{\raggedleft\arraybackslash}p{#1}}

%% Числа: десятичная запятая и выравнивание по разрядам в столбцах S.
\sisetup{
    detect-all,
    output-decimal-marker={,},
    group-separator={\,},
    group-minimum-digits=4,
    table-number-alignment=right
}

%% Продолжение таблицы при переносе longtable на следующую страницу.
\newcommand{\vkrtablecontinued}{\caption*{Продолжение таблицы \thetable}}
